\documentclass[10pt,a4paper]{amsart}
\usepackage[margin=19mm]{geometry}
\usepackage{amsmath,amssymb,mathtools}
\usepackage[scr=rsfs,cal=euler]{mathalfa}
\usepackage{xfrac}
\usepackage[unicode,hidelinks,bookmarks=false]{hyperref}
\newcommand{\ie}{i.e.\ }
\newcommand{\cf}{cf.\ }
\newcommand{\resp}{respectively\ }
\newcommand{\Subsection}{\subsection}
\providecommand{\nobreakdash}{\nobreak\hbox{-}}
\setlength{\emergencystretch}{2em}
\multlinegap0pt
\usepackage{bm}
\usepackage{bbm}
\usepackage{dsfont}
\usepackage{xspace}
\usepackage{cancel}
\usepackage{mathtools}
\usepackage[shortlabels]{enumitem}
\usepackage{amsthm}
\makeatletter
\@ifundefined{definition}{\newtheorem{definition}{Definition}}{}
\@ifundefined{lemma}{\newtheorem{lemma}[definition]{Lemma}}{}
\@ifundefined{claim}{\newtheorem{claim}[definition]{Claim}}{}
\makeatother
\makeatletter
\newcommand{\dist}{\operatorname{dist}}
\expandafter\ifx\csname poc\endcsname\relax
\newenvironment{poc}{\begin{proof}[Proof of the claim]\renewcommand*{\qedsymbol}{$\blacksquare$}}{\end{proof}}
\fi
\makeatother
\title[On cliques in hypergraphs]{On cliques in hypergraphs}
\author{Jun Gao}
\date{Source: arXiv:2510.14804v3 (CC BY 4.0)}
\begin{document}
\maketitle

\noindent\emph{Source and licence.} On cliques in hypergraphs. Version arXiv:2510.14804v3 (CC BY 4.0), distributed under the Creative Commons Attribution 4.0 International licence, as recorded in the arXiv licence metadata for this exact version and shown on the abstract page https://arxiv.org/abs/2510.14804v3. Full rights notice: licenses/arxiv-2510.14804-CC-BY-4.0.txt.\par\smallskip

\noindent\textit{Editorial scope.} Counted unit: the Lemma whose source label is \texttt{lem: key} in the article (source0.tex lines 221--223, source label \texttt{lem: key}), delivered together with the article's own complete original proof of it (source0.tex lines 224--277, one \texttt{proof} environment of 54 lines). No mathematics is written, repaired or re-derived: statement and proof are the article's own text line for line. Required context retained uncounted: def: tree (lines 211--218). Every definition, display and label of those lines is retained unchanged. Omitted from this package and not redistributed: 1--210, 219--220, 278--437. Nothing inside a retained excerpt is omitted or altered: every retained body is delivered exactly as the article's lines are written, so no omission receipt of any kind accompanies this package. Citations: the retained excerpts carry no citation command at all, so no bibliography entry of the article is transcribed and no reference is redistributed. Reference adaptation: none was needed and none was made; every reference inside the delivered unit resolves inside it, so no unresolved reference or citation is produced. Printed slips kept literal: no printed word, symbol, hypothesis or number of the counted statement or of its proof was altered. Layout and class: the article's own class is \texttt{article} with packages including epsf, epsfig, amsfonts, amsgen, amsmath, amstext, amsbsy, amsopn, amsthm, cases, listings, color, ebezier, eepic; the wrapper is the lane's pinned \texttt{amsart} editorial wrapper at 10pt on A4, which restates the theorem-like environments the retained text needs and adds the article's own macro definitions that the retained text needs (\texttt{\textbackslash dist}), with extra wrapper packages (bm, bbm, dsfont, xspace, cancel, mathtools, enumitem). The delivered numbering is the unit's own. Assets: no retained span contains a figure environment, a graphics-inclusion command, a PDF-page inclusion, a table or a TikZ picture; no image file is copied into this package, the package declares no asset dependency and no image file is redistributed with it. Licence: version arXiv:2510.14804v3 of this article is distributed under the Creative Commons Attribution 4.0 International licence, as recorded in the arXiv licence metadata for this exact version and shown on the abstract page https://arxiv.org/abs/2510.14804v3. A full rights notice is delivered at licenses/arxiv-2510.14804-CC-BY-4.0.txt. Characters: every retained excerpt is pure ASCII, so the delivered engine, XeLaTeX with its default Latin Modern text font, reports no missing character.
\begin{definition}[$(k,C)$-layered tree]\label{def: tree}
    We say a tree $T$ is a $(k,C)$-layered tree if we can order vertices of the tree with $v_0,v_1,v_2,\cdots,v_t$ such that the following holds:
    \begin{enumerate}[label=(\arabic*)]
    \item Each vertex $v_i$ is a neighbor of some $v_j$ with $j<i$.
    \item For any $i\in [t]$, $\dist(v_0,v_i) \le k$.
    \item For each vertex $v_i\in  V(T)$, the degree of $v_i$ in $T$ is at most $2^{C+i}$.
    \end{enumerate}
\end{definition}

\begin{lemma}\label{lem: key}
    There exists a constant $N_0=N_0(C,k)$ such that any $(k,C)$-layered tree has at most $N_0$ vertices.
\end{lemma}

\begin{proof}
    Let $T$ be a $(k,C)$-layered tree with the vertices $v_0,v_1,v_2,\cdots,v_t$ satisfying the condition in Definition \ref{def: tree}.
    The proof proceeds by induction on $k$.
    
    When $k=1$, for any $i\in[t]$, $v_i$ is the neighbor of $v_0$ as $\dist(v_0,v_i) \le 1$. 
    It follows from $\deg_T(v_0) \le 2^C$ that
    $|V(T)| = t+1 \le 2^C+1$, so Lemma~\ref{lem: key} holds for $k=1$.

    Assume Lemma~\ref{lem: key} holds for $k-1$. Let $v_{x_1},v_{x_2}, \dots, v_{x_m}$ be the neighbors of $v_0$ in $T$ such that $x_1 < x_2<\dots < x_m$.
    Set $x_{m+1} = t+1$.
    For each $i\in[m]$.
    Let $T_i$ be an subtree of  $T$ with vertex set $\{v_0,v_1,v_2,\dots, v_{x_{i+1}-2}, v_{x_{i+1}-1}\}$. Let $T'_i$ be the tree obtained from $T_i$ by merging vertices $v_{x_1},v_{x_2}, \dots, v_{x_i}$ into a single new vertex $w_{i,0}$, replacing the edge with one endpoint in $\{v_{x_1},v_{x_2}, \dots, v_{x_i}\}$ by an edge incident to $w_{i,0}$ and then deleting the vertex $v_0$. Then we have $|V(T'_i)| = x_{i+1}-i$.
    \begin{claim}\label{claim: 2.3}
        If $\deg_T(v_{x_j}) \le C_j$ for any $j\in [i]$, then $T'_i$ is a $(k-1,2^C + C + \sum^{i}_{j=1}C_j)$-layered tree.
    \end{claim}
    \begin{poc}
    Given an integer $i$, we set $w_{i,0} = w_0$.
    The remaining vertices in the tree $T'_i$ are ordered in such a way $w_1,w_2,\cdots,w_{x_{i+1}-i-1}$ that their relative order is the same as in the original sequence, i.e., if $w_{i_1} =v_{j_1}$ and $w_{i_2} = v_{j_{2}}$, then $i_1<i_2$ if and only if $j_1<j_2$.
    Next, we prove that $T'_i$ is a $(k-1,2^C + C + \sum^{i}_{j=1}C_j)$-layered tree by verifying the three conditions in the definition with ordering $w_0,w_1,w_2\cdots,w_{x_{i+1}-i-1}$ .
    The first two conditions are automatically satisfied, as we do not change the order of the vertices except for $w_0$ and the vertex $v_0$ is removed. 
    
    Note that
        \[
        \deg_{T'_i}(w_0) = \sum_{j=1}^{i} \left(\deg_{T_i}(v_{x_j}) -1 \right) \le \sum_{j=1}^{i} \left(\deg_{T}(v_{x_j}) -1 \right) < \sum^{i}_{j=1}C_j\le 2^{\sum^{i}_{j=1}C_j}, 
        \]
        thus the third condition holds for $w_0$.
        For each $w_j \ne w_0$, there exists $j' \le j+2^C$ such that $w_j =v_{j'}$, as the size of the vertex set we are contracting is $i$, which satisfies $i\le m \le 2^C$.
        Note that for each $w_j \in T'_i$ with $w_j\ne w_0$, the degree of $w_j$ in $T'_i$, is at most the degree of $w_j=v_{j'}$ in $T$.
        We derive that
        \[
        \deg_{T'_i}(w_j)\le \deg_{T}(v_{j'})\le 2^{j'+C} \le 2^{j+2^C+C} \le  2^{j+2^C + C + \sum^{i}_{j=1}C_j}.
        \]
        So $T'_i$ is a $(k-1,2^C + C + \sum^{i}_{j=1}C_j)$-layered tree.
    \end{poc}
    \begin{claim}\label{claim: 2.4}
        For each $j\in [m]$, there exists a constant $C_j=C_j(C,k)$ such that $\deg_T(v_{x_j}) \le C_j$.
    \end{claim}
    \begin{poc}
        Let $m'$ be the maximum integer such that there exists a constant $C_{m'}=C_{m'}(C,k)$ such that $\deg_T(v_{x_{m'}}) \le C_{m'}$. It follows from $v_1=v_{x_1}$ is a neighbor of $v_0$ that we can take $C_1 = 2^{C+1}$, which implies that such a $m'$ exists.

        Suppose to the contrary that $m' <m$. By Claim~\ref{claim: 2.3}, $T'_{m'}$ is a $(k-1,2^C + C + \sum^{m'}_{j=1}C_j)$-layered tree.
        By the inductive hypothesis, we know that there exists a constant \[N_{m'} = N_0\left(k-1,2^{2^C + C + \sum^{m'}_{i=1}C_i}\right) \text{ such that }
        |V(T'_{m'})| = x_{m'+1}-m' \le N_{m'}.\] 
        Take $C_{m'+1} = 2^{N_{m'} +2^C+C} $, then $C_{m'+1}$ is a constant in terms of $k$ and $C$ with
        \[\deg_{T}(v_{x_{m'+1}}) \le 2^{x_{m'+1}+C} \le 2^{N_{m'}+m'+C} \le C_{m'+1}, \]
        a contradiction.
    \end{poc}
    By Claim~\ref{claim: 2.3}, $T'_m$ is a $(k-1,2^C + C + \sum^{m}_{i=1}C_i)$ layered tree. 
    By the inductive hypothesis and Claim~\ref{claim: 2.4}, there exists a constant $N_m$ in terms of $k$ and $C$ such that $|V(T'_m)| = t+1-m \le N_m$, which implies that
    \[
    |V(T)| = t+1 \le N_m+m \le N_m+2^C.
    \]
    So Lemma~\ref{lem: key} holds for $k$.
\end{proof}

\end{document}
